\section{Binary Responses}
\label{app:binary}
\label{sec:binary}

The transition from a factor favoring one trait to a factor representing both traits also appeared in population computations for binary items. We considered four normal-ogive items, two per trait, with slopes corresponding to a loading of .7. The population distribution was evaluated over all $2^p$ response patterns, and a one-trait model was fitted by minimizing KL divergence with multiple starts and 41 to 121 Gauss--Hermite nodes; the fits reported below are the best solutions found. The ratio $1-\KL/\TC$, termed information CFI, extends the population CFI interpretation to this response format.

With thresholds fixed at 0 and $\varphi=0$, the best-found fit had slopes 0 in one block and .98 in the other. The information loss was 0.0781 bits, equal to the mutual information within the ignored block, at each of 41, 81 and 121 quadrature nodes. At $\varphi=.2$, .3 and .35, the favored and other block slopes were .94 and .23, .87 and .37, and .80 and .47, respectively. At $\varphi=.4$, the slopes were equal. These computations place the binary transition between .35 and .40, below the Gaussian threshold of .49 for the corresponding design; they do not establish a general binary threshold theorem.

The information loss was almost entirely residual local dependence. Coupling the true response distribution with the fitted model's posterior distribution of the trait decomposes KL divergence into residual local dependence given the trait, item-response-function misfit, and latent-distribution misfit. The three components summed to the KL divergence within floating-point rounding (absolute difference below $10^{-15}$). Residual local dependence accounted for more than 99.5\% of the one-trait discrepancy at every evaluated $\varphi$, both with thresholds fixed at 0 and in the design of Table~\ref{tab:binary-information}.

\begin{table}[t]
\centering\small
\caption{Information retained by a one-trait normal-ogive model for four binary items generated by two correlated traits. The generating loading is .7, so the normal-ogive slope is $a=.98$; the item intercepts are $b=-.5$, .5, $-.3$ and .3 in $P(X=1\mid\theta)=\Phi(a\theta-b)$, corresponding to standardized latent-response thresholds of $-.36$, .36, $-.21$ and .21. Information is reported in bits, and the one-trait fits are the best solutions found. Sample sizes, rounded to three significant digits, use a noncentral-$\chi^2$ approximation to the full-information $G^2$ test at $\alpha=.05$, with seven degrees of freedom; finite-sample accuracy of this approximation was not assessed.}
\label{tab:binary-information}
\begin{tabular}{cccccc}
\toprule
$\varphi$ & $H(X)$, bits & $\TC(X)$, bits & KL lost by one trait, bits & Information CFI & \shortstack{Approximate $N$\\for 80\% power} \\
\midrule
.30 & 3.689 & 0.156 & 0.0545 & .651 & 190 \\
.50 & 3.663 & 0.182 & 0.0335 & .816 & 310 \\
.70 & 3.621 & 0.224 & 0.0139 & .938 & 750 \\
.80 & 3.593 & 0.251 & 0.00671 & .973 & 1,550 \\
.90 & 3.560 & 0.284 & 0.00183 & .994 & 5,660 \\
.95 & 3.541 & 0.303 & 0.000479 & .998 & 21,600 \\
\bottomrule
\end{tabular}
\end{table}

At high trait correlations, little response information distinguished the two-trait model from its one-trait approximation. Table~\ref{tab:binary-information} uses the same loading design with item intercepts $-.5$, .5, $-.3$ and .3. The response entropy was about 3.6 of the 4 possible bits. Inter-item dependence accounted for 0.16 to 0.30 bits, and the additional information lost by fitting one trait was about 0.05 bits or less. The approximate sample size for 80\% power consequently increased sharply with the trait correlation. These sample sizes use a noncentral-$\chi^2$ approximation to the full-information $G^2$ test, with seven degrees of freedom, $\alpha=.05$, and noncentrality $2N\cdot\KL$ in nats. Their finite-sample accuracy, including behavior near transition points or nonregular endpoints, was not assessed.
